\section{CMSIS-FreeRTOS v2}
System FreeRTOS to popularny system operacyjny czasu rzeczywistego o otwartym kodzie źródłowym dedykowanego dla systemów wbudowanych i mikrokontrolerów. Umożliwia zarządzanie wieloma zadaniami jednocześnie. FreeRTOS działa w oparciu o zadaniowy model programowania, w którym każde zadanie utworzone w programie może mieć przypisany priorytet. To pozwala na efektywne zarządzanie zasobami systemu i nadawanie priorytetu krytycznym operacjom. \newline 
\noindent W centrum FreeRTOS znajduje się harmonogram odpowiadający za decydowanie wykonywaniem konkretnych zadań. Dzięki temu FreeRTOS może być elastycznie dostosowywany do różnych aplikacji.
System oferuje także zaawansowane mechanizmy zarządzania zasobami, takie jak semafory, czy muteksy. Za ich pomocą możemy zezwolić na bezpieczną komunikację między zadaniami i synchronizują ich działanie. Podobnie w przypadku kolejek umożliwiają one przesyłanie danych, co jest niezbędne w aplikacjach wymagających koordynacji różnych procesów. \newline 
\noindent FreeRTOS jest niezbędny do obsługi serwera na STM32, ponieważ zapewnia wielozadaniowość, deterministyczne zarządzanie zadaniami oraz efektywne wykorzystanie zasobów mikrokontrolera.
